\section{Parámetros, unidades y procedencia}

\subsection{Parámetros del núcleo larvario}

\begin{table}[H]
\centering
\small
\caption{Parámetros principales del núcleo de Eriksen y su tratamiento
en la calibración propia.}
\begin{tabularx}{\textwidth}{l l X l}
\toprule
Símbolo & Unidades & Significado & Tratamiento \\
\midrule
$B$ & mg C/larva & Biomasa estructural & Estado \\
$L_{\mathrm{tot}}$ & mg C/larva & Lípidos totales & Estado \\
$a_{\max}$ & d$^{-1}$ & Asimilación máxima de referencia &
Calibrado: $1{,}4995$ \\
$\alpha,\beta$ & -- & Factores de las funciones de crecimiento y
desarrollo & Heredados \\
$Y_B$ & -- & Eficiencia de conversión a biomasa estructural &
Fijo: $0{,}44$ \\
$Y_L$ & -- & Eficiencia de conversión a lípidos &
Fijo: $0{,}42$ \\
$m$ & d$^{-1}$ & Coeficiente de mantenimiento &
Fijo: $0{,}08$ \\
$B_{\max,0}$ & mg C & Escala de biomasa estructural máxima &
Heredado: $39{,}69$ \\
$t_{p,\min}$ & d & Umbral mínimo de desarrollo &
$8{,}5$ desde montaje \\
$\rho$ & mg C/d & Escala del costo máximo de crecimiento &
Calibrado: $4{,}395$ \\
$\delta_I$ & d & Duración de instares & Heredado \\
$\delta_{CL}$ & -- & Fracción de carbono en lípidos &
Heredado: $0{,}79$ \\
$\delta_{CB}$ & -- & Fracción de carbono en biomasa &
Heredado: $0{,}49$ \\
\bottomrule
\end{tabularx}
\end{table}

La conservación del núcleo no implica conservar idénticos todos sus
valores numéricos. En la calibración propia se ajustan $a_{\max}$,
$\rho$ y $f_{D4}$, mientras $Y_B$, $Y_L$ y $m$ permanecen congelados.
Por tanto, no es correcto afirmar que el modelo propuesto eliminó los
costos de crecimiento o el mantenimiento.

Esta distinción también explica parte de las diferencias externas. En
la reproducción de Laganaro, Eriksen emplea $Y_B=0{,}44$ y $m=0{,}08$
para la dieta de alta calidad, pero $Y_B=0{,}76$ y $m=0{,}23$ para el
escenario con 75\,\% de lodos. Esos cambios aumentan los costos
metabólicos asociados con la dieta de baja calidad. El modelo propuesto
actual, con $Y_B$ y $m$ congelados, no reproduce ese mecanismo.

\subsection{Parámetros del alimento y del lecho}

\begin{table}[H]
\centering
\small
\caption{Parámetros de la extensión alimento--lecho.}
\begin{tabularx}{\textwidth}{l l X l}
\toprule
Símbolo & Unidades & Significado & Valor u origen \\
\midrule
$D_C$ & g C/cámara & Carbono disponible del alimento & Estado \\
$X$ & g C/cámara & Biomasa microbiana & Estado \\
$f_d$ & -- & Disponibilidad integral de la dieta &
D1: $1$; D4: $0{,}8047$ \\
$K_D$ & g C & Semisaturación del acople alimento--larva &
Fijo: $2$ en tratamientos \\
$k_{\mathrm{ref}}$ & d$^{-1}$ & Cinética microbiana de referencia &
D1: $0{,}6086$; D4: $1{,}0$ \\
$K_{SM}$ & g C & Semisaturación de degradación & Fijo: $50$ \\
$Y_X$ & -- & Rendimiento de biomasa microbiana & Fijo: $0{,}4$ \\
$Y_{C,\mathrm{mic}}$ & -- & Fracción degradada emitida como \co{} &
Fijo: $0{,}4$ \\
$Q_{10}$ & -- & Factor de dependencia térmica & Fijo: $2$ \\
$T_{\mathrm{ref}}$ & $^\circ$C & Temperatura de referencia &
Fijo: $25$ \\
$\theta_d$ & -- & Humedad efectiva adimensional &
Según escenario \\
$K_\theta$ & -- & Semisaturación por humedad & Fijo: $0{,}15$ \\
$\delta_{CS}$ & -- & Fracción de carbono disponible del sustrato &
D1: $0{,}95$; D4: $0{,}2483$ \\
\bottomrule
\end{tabularx}
\end{table}

El valor $a_{\max}=1{,}4995$ d$^{-1}$ se encuentra próximo a su cota
superior. También $k_{\mathrm{ref}}(D4)=1$ y
$\delta_{CS}(D1)=0{,}95$ alcanzan sus fronteras superiores. Estos
resultados se reportan como tales: no se eliminan ni se sustituyen por
valores más favorables. Señalan que algunos parámetros absorben
limitaciones estructurales o información insuficiente.

\subsection{Qué se calibró en cada etapa}

La primera etapa utiliza controles sin larvas para estimar
$k_{\mathrm{ref}}$ y $\delta_{CS}$ por dieta. La segunda congela esos
parámetros microbianos y ajusta $a_{\max}$, $\rho$ y $f_{D4}$ contra
crecimiento inicial y ocho cierres de gas. El ajuste de crecimiento se
restringe a días $\leq8$; la fase tardía se conserva como diagnóstico.

La mejora obtenida en la segunda etapa corresponde al sistema completo.
No identifica por sí sola la contribución causal del lecho, porque la
recalibración larvaria y el acople con el alimento cambian
simultáneamente.
